%% SCRAP: papers/NEGATIVE_RESULTS
%% SOURCE: docs/working/papers/NEGATIVE_RESULTS.md
%% STATUS: CURRENT
%% FITS: ssrn/ch-claims, vol3-research/ch-formal-claims
%% EDITORIAL: lifted — prose rewritten to press voice

\section{Negative Results: Documented Failure Modes}
\label{sec:negative-results}

Systems that never fail are artifacts, not scientific results. This section
documents 15 failure modes identified through deliberate adversarial testing.
These are not accidental discoveries; experiments were designed to break the
system. The failure modes bound the system's applicability and constitute
essential context for interpreting the positive claims.

A concise summary table appears at the end of this section
(\S\ref{sec:failure-summary}).

\subsection{Workload-Dependent Failures}

\paragraph{Failure mode 1: random execution paths.}
A workload in which control flow depends on a nanosecond timer
(\texttt{TIMER @ 2 MOD}) produces non-reproducible execution sequences.
Cache hit rate CV rises to approximately 15\%; execution heat distributions
are non-reproducible; the system oscillates without converging. Root cause:
the rolling window captures different execution sequences on each run, giving
the adaptive mechanisms a moving target. No mitigation exists; this is an
intentional scope boundary. Deterministic adaptation requires deterministic
workloads.

\paragraph{Failure mode 2: I/O-bound workloads.}
A workload dominated by file I/O (1{,}000 iterations of
\texttt{OPEN-FILE~/ READ-FILE / CLOSE-FILE}) achieves at most 2\%
improvement, within margin of error. The adaptive mechanisms add approximately
5\% overhead, producing net degradation. Root cause: the bottleneck is disk
latency, not dictionary lookup; I/O timing variance swamps algorithmic
determinism. Mitigation (not yet implemented): detect I/O-bound workloads
and disable adaptive loops.

\paragraph{Failure mode 3: short programs.}
Programs with fewer than approximately 1{,}000 word executions provide
insufficient data for statistical inference. Levene's test requires minimum
sample sizes; exponential regression on three data points is meaningless.
Overhead exceeds benefit. Mitigation: disable adaptive loops when
\texttt{execution\_count~<~THRESHOLD} (suggested: 10{,}000).

\paragraph{Failure mode 4: adversarial workloads.}
A workload that executes each word exactly once in rotation produces a flat
frequency distribution---the Zipf-law assumption is violated. The hot-words
cache is useless (all words equally ``hot''); performance improvement is 0\%.
Mitigation: detect flat distributions via entropy measurement and disable the
hot-words cache.

\subsection{Parameter-Dependent Failures}

\paragraph{Failure mode 5: window size too small.}
Setting \texttt{ROLLING\_WINDOW\_SIZE = 10} (default: 4{,}096) produces
unstable variance inflection detection ($R^2 < 0.5$) and convergence
oscillation. Statistical noise dominates the signal. Mitigation: enforce a
minimum window size of at least 1{,}024 entries.

\paragraph{Failure mode 6: decay slope too steep.}
Setting $\lambda = 10.0$ (default: ${\approx}0.001$) causes cache thrashing:
words are promoted and immediately demoted, the system never reaches steady
state, and performance degrades below baseline. Root cause: steep decay erases
long-term frequency information, making the cache reactive to noise rather
than signal. Mitigation: bound decay slope to $\lambda \in [10^{-4}, 10^{-2}]$.

\paragraph{Failure mode 7: cache size mismatch.}
Setting \texttt{HOTWORDS\_CACHE\_SIZE = 1} (default: 16) reduces cache hit
rate to approximately 5\% and performance improvement to approximately 2\%.
Root cause: the test workload contains 10--15 hot words; a cache of 1 entry
captures only the single hottest. Mitigation: auto-tune cache size based on
workload entropy.

\subsection{Environmental Failures}

\paragraph{Failure mode 8: thermal throttling.}
Sustained CPU load sufficient to trigger thermal throttling increases runtime
CV to 90\raisebox{0.5ex}{+}\% and extends the convergence window from
${\approx}30$ to ${\approx}40$ runs. Cache decisions remain at 0\% CV:
timing does not affect algorithmic choices. Root cause: wall-clock-based
decay becomes inconsistent when the CPU reduces its clock frequency.
Mitigation: use CPU cycle counters instead of wall-clock timers for
decay application.

\paragraph{Failure mode 9: aggressive OS preemption.}
Running with 32 concurrent CPU-saturating background processes pushes runtime
CV above 150\% and may prevent convergence if the process is preempted too
frequently for decay to apply correctly. Cache decisions remain at 0\% CV.
Mitigation: pin the process to an isolated core and set real-time priority.

\subsection{Architectural Failures}

\paragraph{Failure mode 10: cross-architecture performance difference.}
Cache decisions are bit-identical across x86\_64 (Intel Xeon) and AArch64
(ARM Cortex) because the algorithm is purely arithmetic. Convergence
\emph{rate} differs---Intel achieves approximately 25\% improvement, ARM
approximately 18\%---because dictionary lookup has a different relative cost
on each architecture. This is expected behavior, not a failure of
determinism. No mitigation is needed; cross-architecture performance
comparisons require hardware-normalized baselines.

\paragraph{Failure mode 11: floating-point non-determinism (hypothetical).}
IEEE-754 arithmetic is not used anywhere in the adaptive runtime; \Qtype\
fixed-point arithmetic provides bit-identical results across architectures
and compilers. This entry documents why floating-point was avoided: FMA
instructions, compiler reordering, and denormal handling all introduce
non-determinism that would violate the 0\% CV claim.

\subsection{Implementation Bugs Resolved}

\paragraph{Bug 1: heartbeat thread race condition (fixed).}
An unsynchronized shared-state access between the heartbeat thread and the
main interpreter thread caused segmentation faults in fewer than 1\% of
runs. Fix: added a mutex around \texttt{vm\_tick()} (commit
\texttt{8133787}). The experimental data was collected after this fix.

\paragraph{Bug 2: integer overflow in execution heat (fixed).}
32-bit frequency counters overflowed on long-running programs, causing
execution frequency to reset to zero unexpectedly. Fix: counters promoted
to \texttt{uint64\_t} (commit \texttt{c0ec82d}).

\subsection{Statistical Failures}

\paragraph{Failure mode 12: insufficient sample size.}
With $N = 5$ trials instead of the $N = 30$ used in the experiments, the
Central Limit Theorem does not apply; confidence intervals are unreliable;
$t$-test normality assumptions are violated; statistical power falls below
50\%. Mitigation: require $n \geq 30$ for all experiments.

\paragraph{Failure mode 13: multiple testing without correction.}
Testing 100 hypotheses at $\alpha = 0.05$ without Bonferroni correction
yields approximately 5 spurious positives by chance. The formal claim table
tests 5 primary claims; corrected $\alpha = 0.05 / 5 = 0.01$. All reported
$p$-values remain significant under this correction.

\subsection{Generalization Failures}

\paragraph{Failure mode 14: compiled languages.}
The adaptive runtime is an interpreter optimization. Ahead-of-time compilers
already perform static frequency analysis and hot-code optimization without
runtime overhead; the rolling window and heartbeat infrastructure would add
cost with no benefit. This is a scope boundary, not a bug.

\paragraph{Failure mode 15: very large programs.}
Programs with 100{,}000\raisebox{0.5ex}{+} dictionary entries would require
sorting a 100K-entry array on every heartbeat tick and maintaining a
$100\text{K} \times 100\text{K}$ transition matrix---both computationally
prohibitive. Scalability at this scale is unvalidated; it is marked as
future work.

\subsection{Summary of Failure Modes}
\label{sec:failure-summary}

\begin{table}[h]
\centering
\caption{Failure mode summary. All failures are predictable and bounded.}
\label{tab:failure-modes}
\begin{tabular}{lll}
\toprule
\textbf{Mode} & \textbf{Root cause} & \textbf{Mitigation} \\
\midrule
Random execution paths & Non-deterministic workload & Detect and disable \\
I/O-bound program & Wrong bottleneck & Profile first \\
Short program & Insufficient data & Minimum execution threshold \\
Adversarial (flat) workload & No frequency skew & Detect and disable cache \\
Window size too small & Insufficient context & Enforce minimum: 1{,}024 \\
Decay too steep & Premature forgetting & Bound $\lambda \in [10^{-4}, 10^{-2}]$ \\
Cache too small & Mismatched capacity & Auto-tune to workload entropy \\
Thermal throttling & Inconsistent clock & Use cycle counters \\
OS preemption & Context switches & Isolate core, RT priority \\
Cross-architecture & Hardware differences & Expected; document \\
Floating-point (hypothetical) & IEEE-754 rounding & Fixed-point throughout \\
Small sample ($N < 30$) & CLT does not apply & Require $n \geq 30$ \\
Multiple testing & Inflated $\alpha$ & Bonferroni correction \\
Compiled languages & Wrong paradigm & Out of scope \\
Massive programs & Untested regime & Future work \\
\bottomrule
\end{tabular}
\end{table}

\subsection{Lessons}

Five lessons emerge from the failure modes:

\begin{itemize}
  \item Determinism requires deterministic inputs: random workloads break all
    adaptive guarantees.
  \item Statistical inference requires sufficient data: short programs and
    small sample sizes produce meaningless results.
  \item Optimization requires locality: workloads with flat frequency
    distributions have nothing to cache.
  \item Architecture matters: floating-point and thermal effects are real
    engineering constraints, not theoretical concerns.
  \item Scope boundaries make positive results credible: a system claiming
    unlimited applicability invites justified skepticism.
\end{itemize}

The claims in \S\ref{sec:formal-claims} are scoped to CPU-bound, deterministic,
long-running FORTH programs. Outside that scope, the system fails
predictably and in ways documented above.
